\grechangedim{abovelinestextraise}{0.2cm}{scalable}
\gresetinitiallines{0}

\subsection{Primeira família: tom de Dó}

Corda de recitação dos versos e nota final do responsório em Dó (ou Fá)

\subsubsection{C 1}

\begin{rubrica}
  (Para salmos responsoriais e aleluiáticos com no máximo um asterisco por estrofe)
\end{rubrica}

\gabcsnippet{(c4) <c><sp>R/</sp>.</c>() (hrhrhr/[10][alt:Refrão longo]) (g) (f) (g) (hr) (h) (,) (hrhrhr/[10][alt:Refrão breve]) (g) (f) (fr) (er) (d) (cr) (f) (g) (g) (fr) (f) <c>ou</c>(::) A(h)le(hg)lu(g>)ia,(f) (,) a(d)le(fg)lu(g>)ia.(f) (::)}

\gabcsnippet{(c4) <c><sp>V/</sp>.</c>() <c><v>(</v></c>() (frfrfr/[10]) (fr1) (fr) (fr) +<c><v>)</v></c>(,) (frfrfr/[10]) (gr1) (fr) (f) *(:) (frfrfr/[10]) (e) (gr1) (gr) (h) (::)}

\subsubsection{C 2 a}

\begin{rubrica}
  (Para salmos responsoriais com no mínimo um asterisco por estrofe)
\end{rubrica}

\gabcsnippet{(c4) <c><sp>R/</sp>.</c>() (hrhrhr/[10][alt:Refrão longo]) (g) (f) (g) (hr) (h) (,) (hrhrhr/[10][alt:Refrão breve]) (g) (f) (fr) (er) (d) (cr) (f) (g) (g) (fr) (f) (::)}

\gabcsnippet{(c4) <c><sp>V/</sp>.</c>() <c><v>(</v></c>() (frfrfr/[10]) (fr1) (fr) (fr) +<c><v>)</v></c>(,) (frfrfr/[10]) (gr1) (fr) (f) *(:) (frfrfr/[10]) (gr[ocba:1{]) (gh[ocba:0}]) (h) (::)}

\subsubsection{C 2 g}

\begin{rubrica}
  (Para salmos responsoriais com no mínimo um asterisco por estrofe)
\end{rubrica}

\gabcsnippet{(c4) <c><sp>R/</sp>.</c>() (frfrfr/[10][alt:Refrão longo]) (g) (hr) (h) (,) (hrhrhr/[10]) (g) (f) (fr) (er) (d) (cr) (f) (g) (g) (fr) (f) <c>ou</c>(::) (frfrfr/[10][alt:Refrão breve]) (g) (h) (g) (f) (fr) (er) (d) (cr) (f) (g) (g) (fr) (f) (::)}

\gabcsnippet{(c4) <c><sp>V/</sp>.</c>() <c><v>(</v></c>() (frfrfr/[10]) (fr1) (fr) (fr) +<c><v>)</v></c>(,) (frfrfr/[10]) (gr1) (fr) (f) *(:) (frfrfr/[10]) (gr[ocba:1{]) (gh[ocba:0}]) (hg) (::)}

\subsubsection{C 3 a}

\begin{rubrica}
  (Para salmos responsoriais com no mínimo dois asteriscos por estrofe)
\end{rubrica}

\gabcsnippet{(c4) <c><sp>R/</sp>.</c>() (hrhrhr/[10][alt:Refrão longo]) (g) (f) (g) (hr) (h) (,) (hrhrhr/[10][alt:Refrão breve]) (g) (f) (fr) (er) (d) (cr) (f) (g) (g) (fr) (f) (::)}

\gabcsnippet{(c4) <c><sp>V/</sp>.</c>() <c><v>(</v></c>() (frfrfr/[10][alt:Antepenúltimo e penúltimo asterisco]) (fr1) (fr) (fr) +<c><v>)</v></c>(,) (frfrfr/[10]) (gr1) (fr) (f) *(:) (frfrfr/[10]) (gr[ocba:1{]) (gh[ocba:0}]) (h) (::) (hr[alt:Último asterisco]) (ixhr1i) <c><v>(</v></c>() (hrhrhr/[10]) (hr1) (hr) (hr) +<c><v>)</v></c>(,) (hrhrhr/[10]) (gr1h) (gr) (g) *(:) (grgrgr/[10]) (f) (gr[ocba:1{]) (gh[ocba:0}]) (h) (::)}

\subsubsection{C 3 g}

\begin{rubrica}
  (Para salmos responsoriais com no mínimo dois asteriscos por estrofe)
\end{rubrica}

\gabcsnippet{(c4) <c><sp>R/</sp>.</c>() (frfrfr/[10][alt:Refrão longo]) (g) (hr) (h) (,) (hrhrhr/[10]) (g) (f) (fr) (er) (d) (cr) (f) (g) (g) (fr) (f) <c>ou</c>(::) (frfrfr/[10][alt:Refrão breve]) (g) (h) (g) (f) (fr) (er) (d) (cr) (f) (g) (g) (fr) (f) (::)}

\gabcsnippet{(c4) <c><sp>V/</sp>.</c>() <c><v>(</v></c>() (frfrfr/[10][alt:Antepenúltimo asterisco]) (fr1) (fr) (fr) +<c><v>)</v></c>(,) (frfrfr/[10]) (gr1) (fr) (f) *(:) (frfrfr/[10]) (gr[ocba:1{]) (gh[ocba:0}]) (hg) (::) <c><v>(</v></c>() (frfrfr/[10][alt:Penúltimo asterisco]) (fr1) (fr) (fr) +<c><v>)</v></c>(,) (frfrfr/[10]) (gr1) (fr) (f) *(:) (frfrfr/[10]) (gr[ocba:1{]) (gh[ocba:0}]) (h) (::) (hr[alt:Último asterisco]) (ixhr1i) <c><v>(</v></c>() (hrhrhr/[10]) (hr1) (hr) (hr) +<c><v>)</v></c>(,) (hrhrhr/[10]) (gr1h) (gr) (g) *(:) (grgrgr/[10]) (f) (gr[ocba:1{]) (gh[ocba:0}]) (hg) (::)}

\subsubsection{C 4}

\begin{rubrica}
  (Para salmos aleluiáticos com no máximo um asterisco por estrofe)
\end{rubrica}

\gabcsnippet{(c4) <c><sp>V/</sp>.</c> <c><v>(</v></c>() (frfrfr/[10]) (fr1) (fr) (fr) +<c><v>)</v></c>(,) (frfrfr/[10]) (e) (gr1) (gr) (h) (::) <c><sp>R/</sp>.</c> A(f)le(f)lu(d)ia.(c) *(::) <c><v>(</v></c>() (frfrfr/[10]) (fr1) (fr) (fr) +<c><v>)</v></c>(,) (frfrfr/[10]) (e) (gr1) (gr) (h) (::) <c><sp>R/</sp>.</c> A(h)le(hg)lu(g>)ia.(f) (::)}

\subsubsection{C \protect\GreStar}

\begin{rubrica}
  (Para salmos responsoriais com no mínimo um asterisco por estrofe)
\end{rubrica}

\gabcsnippet{(f3) <c><sp>R/</sp>.</c> (hrhrhr/[10]) (g) (f) (f) (g) (h) (i) (j) (i) (,) (iririr/[10]) (g) (h) (i) (g) (h) (fr) (f) (::)}

\gabcsnippet{(f3) <c><sp>V/</sp>.</c> <c><v>(</v></c>() (hrhrhr/[10]) (hr1) (hr) (hr) +<c><v>)</v></c>(,) (hrhrhr/[10]) (g) (f) (hr1) (hr) (h) *(:) (hrhrhr/[10]) (hr1) (fr) (f) (::)}

\subsection{Segunda família: tom de Ré}

Corda de recitação dos versos e nota final do responsório em Ré (ou Lá ou Sol)

\subsubsection{D 1 b}

\begin{rubrica}
  (Para salmos responsoriais e aleluiáticos com no máximo um asterisco por estrofe e refrões breves)
\end{rubrica}

\gabcsnippet{(c4) <c><sp>R/</sp>.</c>() (iririr/[10]) (i) (j) (i) (h) (ir) (hr) (g) (h) (i) (i) (hr) (h) (::) <c><sc>ou</sc></c>() A(i)le(j)lu(i)ia,(hg) (,) a(hi)le(i)lu(h)ia.(h) (::)}

\gabcsnippet{(c4) <c><sp>V/</sp>.</c>() <c><v>(</v></c>() (hrhrhr/[10]) (hr1) (hr) (hr) +<c><v>)</v></c>(,) (hrhrhr/[10]) (ir1) (hr) (h) *(:) (hrhrhr/[10]) (gr1) (gr) (i) (::)}

\subsubsection{D 1 e}

\begin{rubrica}
  (Para salmos responsoriais com no mínimo um asterisco por estrofe)
\end{rubrica}

\gabcsnippet{(c4) <c><sp>R/</sp></c>.() (ererer/[10]) (e) (d) (e) (g) (hr) (h) (,) (hrhrhr/[10]) (g) (f) (g) (f) (e) (::)}

\gabcsnippet{(c4) <c><sp>V/</sp></c>.() <c><v>(</v></c>() (hrhrhr/[10]) (hr1) (hr) (hr) +<c><v>)</v></c>(,) (hrhrhr/[10]) (ir1) (hr) (h) *(:) (hrhrhr/[10]) (g) (h) (gr1) (gr) (e) (::)}

\subsubsection{D 1 g}

\begin{rubrica}
  (Para salmos responsoriais com no mínimo um asterisco por estrofe)
\end{rubrica}

\gabcsnippet{(c4) <c><sp>R/</sp>.</c>() (hrhrhr/[10]) (g) (gr) (i) (,) (iririr/[10]) (i) (j) (i) (h) (ir) (hr) (g) (h) (i) (i) (hr) (h) (::)}

\gabcsnippet{(c4) <c><sp>V/</sp></c>.() <c><v>(</v></c>() (hrhrhr/[10]) (hr1) (hr) (hr) +<c><v>)</v></c>(,) (hrhrhr/[10]) (ir1) (hr) (h) *(:) (hrhrhr/[10]) (hr1) (gr) (g) (::)}

\subsubsection{D \protect\GreStar}

\begin{rubrica}
  (Para salmos responsoriais com no mínimo dois asteriscos por estrofe)
\end{rubrica}

\gabcsnippet{(c4) <c><sp>R/</sp></c>.() (grgrgr/[10]) (f) (h) (j) (jr) (i) (g) (,) (iririr/[10]) (j) (h) (i) (h) (g) (gr) (g) (::)}

\gabcsnippet{(c4) <c><sp>V/</sp></c>.() <c><v>(</v></c>() (jrjrjr/[10][alt:Primeiro asterisco]) (jr1) (jr) (jr) +<c><v>)</v></c>(,) (jrjrjr/[10]) (kr1) (jr) (j) *(:) (jrjrjr/[10]) (jr1) (i) (hg) (::) <c><v>(</v></c>() (grgrgr/[10][alt:Segundo e terceiro asterisco]) (gr1) (gr) (gr) +<c><v>)</v></c>(,) (grgrgr/[10]) (hr1) (gr) (g) *(:) (grgrgr/[10]) (hj) (jr1) (i) (hg) (::)}

\subsection{Terceira família: tom de Mi}

Corda de recitação dos versos e nota final do responsório em Mi (ou Si)

\subsubsection{E 1}

\begin{rubrica}
  (Para salmos responsoriais com no máximo um asterisco por estrofe e refrões breves)
\end{rubrica}

\gabcsnippet{(c4) <c><sp>R/</sp></c>.() (grgrgr/[10]) (e) (f) (er) (d) (c) (d) (f) (f) (er) (dr) (er) (e) (::)}

\gabcsnippet{(c4) <c><sp>V/</sp></c>.() <c><v>(</v></c>() (ererer/[10]) (er1) (er) (er) +<c><v>)</v></c>(,) (ererer/[10]) (er1f) (er) (e) *(:) (ererer/[10]) (fd) (er1) (f) (g) (::)}

\subsubsection{E 2 d}

\begin{rubrica}
  (Para salmos responsoriais e aleluiáticos com no mínimo um asterisco por estrofe e refrões breves)
\end{rubrica}

\gabcsnippet{(f3) <c><sp>R/</sp></c>.() (grgrgr/[10]) (f) (i) (h) (g) <c>ou</c>(::) A(g)le(g)lu(g)ia,(g) (,) a(g)le(g)lu(fih)ia.(g) (::)}

\gabcsnippet{(f3) <c><sp>V/</sp></c>.() <c><v>(</v></c> (grgrgr/[10]) (gr1) (gr) (gr) +<c><v>)</v></c>(,) (grgrgr/[10]) (gr1h) (gr) (g) *(:) (grgrgr/[10]) (gr1) (fr) (f) (::)}

%\subsubsection{E 2 e (ou E 5 *)}

%Observação: esta melodia aparece apenas uma vez no Gradual Simples.

%\gabcsnippet{(f3) <c><sp>V/</sp></c>.() (gR/[10]) (gr1h) (gr) (g) *(:) (gR/[10]) (h) (f) (h) (g) (::) <c><sp>R/</sp></c>.() (e) (g) (g) (g) (h) (h) (g) (e) (f) (g) (f) (g) (g) (::)}

\subsubsection{E 3}

\begin{rubrica}
  (Para salmos responsoriais com no mínimo um asterisco por estrofe e refrões breves)
\end{rubrica}

\gabcsnippet{(c4) <c><sp>R/</sp></c>.() (frfrfr/[10]) (f) (g) (h) (ixi) (g) (ixi) (h) (::)}

\gabcsnippet{(c4) <c><sp>V/</sp></c>.() <c><v>(</v></c>() (hrhrhr/[10]) (hr1) (hr) (hr) +<c><v>)</v></c>(,) (hrhrhr/[10]) (hr1) (gr) (g) *(:) (hrhrhr/[10]) (gr1) (fr) (f) (::)}

\subsubsection{E 4}

\begin{rubrica}
  (Para salmos responsoriais e aleluiáticos com no mínimo um asterisco por estrofe)
\end{rubrica}

\gabcsnippet{(c4) <c><sp>R/</sp></c>.() (hrhrhr/[10]) (h) (gr) (hr) (g) (,) (grgrgr/[10]) (f) (g) (h) (ixi) (g) (ixi) (h) <c>ou</c>(::) A(h)le(g)lu(h)ia,(g) (,) a(g)le(f)lu(g)ia,(h) a(ixi)le(g)lu(i)ia.(h) (::)}

\gabcsnippet{(c4) <c><sp>V/</sp></c>.() <c><v>(</v></c>() (hrhrhr/[10]) (hr1) (hr) (hr) +<c><v>)</v></c>(,) (hrhrhr/[10]) (gr1) (fr) (f) *(:) (hrhrhr/[10]) (gr1) (hr) (h) (::)}

\subsubsection{E 5}

\begin{rubrica}
  (Para salmos responsoriais com no mínimo dois asteriscos por estrofe)
\end{rubrica}

\gabcsnippet{(c4) <c><sp>R/</sp></c>.() (ererer/[10]) (f) (d) (e) (f) (g) (,) (grgrgr/[10]) (e) (f) (er) (d) (c) (d) (f) (f) (er) (dr) (er) (e) (::)}

\gabcsnippet{(c4) <c><sp>V/</sp></c>.() <c><v>(</v></c>() (ererer/[10][alt:Antepenúltimo asterisco]) (er1) (er) (er) +<c><v>)</v></c>(,) (ererer/[10]) (er1f) (er) (e) *(:) (ererer/[10]) (er1) (dr) (d) (::) <c><v>(</v></c>() (ererer/[10][alt:Penúltimo asterisco]) (er1) (er) (er) +<c><v>)</v></c>(,) (ererer/[10]) (er1f) (er) (e) *(:) (ererer/[10]) (fd) (er1) (f) (g) (::) (hr) (hr1i) <c><v>(</v></c>() (hr[alt:Último asterisco]) (hr1i) (hrhrhr/[10]) (hr1) (hr) (hr) +<c><v>)</v></c>(,) (hrhrhr/[10]) (gr1h gr) (g) *(:) (grgrgr/[10]) (gf) (er1) (er) (e) (::)}

\subsubsection{E \protect\GreStar}

\begin{rubrica}
  (Para salmos aleluiáticos com no máximo um asterisco por estrofe)
\end{rubrica}

\gabcsnippet{(f3) <c><sp>R/</sp></c>.() A(h)le(g)lu(hi)ia,(hhg) (,) a(f)le(e)lu(fg~)ia.(g) (::)}

\gabcsnippet{(f3) <c><sp>V/</sp></c>.() <c><v>(</v></c>() (iririr/[10]) (ir1) (ir) (ir) +<c><v>)</v></c>(,) (iririr/[10]) (jr1) (ir) (i) *(:) (iririr/[10]) (ih) (gr1) (gr) (g) (::)}

\subsection{Outros tons salmódicos aleluiáticos}

\subsubsection{Tom romano antigo}

\begin{rubrica}
  (Para salmos aleluiáticos com no mínimo dois asteriscos por estrofe. A primeira alternativa de entoação inicial no último segmento é utilizada quando a primeira sílaba tônica da frase ocorrer na segunda posição, enquanto a segunda alternativa é utilizada nos demais casos.)
\end{rubrica}

\gabcsnippet{(c4) <c><sp>R/</sp>.</c>() A(c)le(d)lu(de)ia,(dc) (,) a(e)le(g)lu(gh)ia,(g) a(fe)le(dc)lu(de)ia.(e) (::)}

\gabcsnippet{(c4) <c><sp>V/</sp>.</c>() (c[alt:Antepenúltimo e penúltimo asterisco]) (de) <c><v>(</v></c>() (ererer/[10]) (er1) (er) (er) +<c><v>)</v></c>(,) (ererer/[10]) (de) (dr1) (cr) (c) *(:) (c) (df) (frfrfr/[10]) (f) (d) (fr1) (fr) (e) (::) (dh[alt:Último asterisco]) (ixhr1i) <c>ou</c>(::) (ixdh!iv) <c><v>(</v></c>() (hrhrhr/[10]) (hr1) (hr) (hr) +<c><v>)</v></c>(,) (hrhrhr/[10]) (gr1h) (gr) (g) *(:)
(grgrgr) (gr1h) (grgrgr/[10]) (f) (d) (fr1) (fr) (ed/ffe) (::)}

\subsubsection{Tom moçárabe}

\begin{rubrica}
  (Para salmos aleluiáticos com no máximo um asterisco por estrofe. As três alernativas de entoação inicial dependem da posição da primeira sílaba tônica na frase --- primeira, segunda ou terceira em diante, respectivamente.)
\end{rubrica}

\gabcsnippet{(f3) <c><sp>R/</sp>.</c>() A(f)le(f)lu(e)ia,(c) (,) a(e)le(f)lu(g)ia,(f) (,) a(f)le(fe/ghGF)lu(ef)ia.(f) (::)}

\gabcsnippet{(f3) <c><sp>V/</sp>.</c>() (er1f) <c>ou</c>(::) (f) (fr1ef) <c>ou</c>(::) (fe) (ef) <c><v>(</v></c>() (frfrfr/[10]) (fr1) (fr) (fr) +<c><v>)</v></c>(,) (frfrfr/[10]) (fe) (fg) *(:)
(e) (g) (frfrfr/[10]) (e) (f) (gr1) (gr) (e) (::)}

\subsubsection{Tom arcaico de Ré}

\begin{rubrica}
  (Para salmos aleluiáticos com no mínimo um asterisco por estrofe)
\end{rubrica}

\gabcsnippet{(c3) <c><sp>R/</sp>.</c>() A(e)le(f)lu(hi)ia,(i) (,) a(j)le(k)lu(j)ia,(ih) a(ij)le(j)lu(i)ia.(i) (::)}

\gabcsnippet{(c3) <c><sp>V/</sp>.</c>() (i[alt:Antepenúltimo e penúltimo asterisco]) (hi) <c><v>(</v></c>() (iririr/[10]) (ir1) (hr) (h) +<c><v>)</v></c>(,) (iririr/[10]) (h) (j) (i) (ir1) (fr) (f) *(:) (f) (h) (iririr/[10]) (h) (i) (j) (i) (hr1) (hr) (hi) (::) (i[alt:Último e penúltimo asterisco]) (hi) <c><v>(</v></c>() (iririr/[10]) (ir1) (hr) (h) +<c><v>)</v></c>(,) (iririr/[10]) (h) (j) (i) (ir1) (fr) (f) *(:)
(f) (h) (iririr/[10]) (h) (i) (j) (hr1) (fr) (fg) (::)}

\gresetinitiallines{1}